\documentclass[11pt]{article}
\usepackage[utf8]{inputenc}
\usepackage{amsmath,amssymb,amsthm}
\usepackage[margin=1in]{geometry}
\usepackage{hyperref}
\usepackage{xcolor}

\newcommand{\tideref}[2][]{\href{https://therisensea.org/tides/#2/}{\texttt{#2}}}

\title{Tide retrospective: susceptibilities recover the anharmonic
coefficients}
\author{Tide loop (Claude Fable 5, with GPT-5.6 Sol deliberation)}
\date{2026-08-09}

\begin{document}
\maketitle

\section{Setting}

Every recovery tide so far certified Section~7.4 of the germbij note
(the singular, constructive side). The note's \emph{regular}-case
statement, Theorem~3.1, says the expansions over all observables
determine the Taylor coefficients of the loss at a nondegenerate
minimum by a triangular induction pairing coefficients against
susceptibilities. The seabed's Stage-2 anharmonic programme (from the
primer work in May) already carries the two limits that drive the
induction's first two rungs; this tide reads them backwards.

\section{The thing formalised}

One theorem in
\texttt{Laplace/OneD/RecoveryAnharmonic.lean}\footnote{\texttt{anharmonic\_susceptibility\_recovery}.},
sorry-free with zero warnings. For the anharmonic potential
$\tfrac{\lambda}{2}x^2 + \tfrac{\alpha}{6}x^3 + \tfrac{\gamma}{24}x^4$
(positive $\lambda, \gamma$, discriminant $\alpha^2 < 3\lambda\gamma$),
the seabed knows that the mean susceptibility
$\partial_h \langle x \rangle_h|_{h=0}$ tends to $-1/\lambda$ and the
rescaled covariance susceptibility to $\alpha/\lambda^3$. The theorem:
if these two observable data streams agree eventually for two parameter
triples, then $\lambda_1 = \lambda_2$ and $\alpha_1 = \alpha_2$. This is
the note's Theorem~3.1 recovery, certified at orders two and three: the
data of the expansions pins the jet, coefficient by coefficient. The
quartic coefficient $\gamma$ is the next rung and needs the
fourth-cumulant asymptotic, which the seabed does not yet carry.

\section{Proof strategy}

Entirely the uniqueness-of-limits pattern the thread has now used three
times: transport each limit along the eventual equality
(\texttt{Tendsto.congr'}), identify the limits
(\texttt{tendsto\_nhds\_unique}), and extract the parameters by field
arithmetic ($-1/\lambda$ is injective on positives; then
$\alpha/\lambda^3$ with $\lambda$ pinned). The susceptibility data are
treated as opaque functions of $t$ throughout; nothing about the
underlying Gibbs derivatives is unfolded.

\section{Roadblocks and resolutions}

None: the first first-attempt-clean build of the recovery thread (and
the second of the germbij programme). The consult confirmed the
statement packaging (two separate eventual-equality hypotheses rather
than a pair) before the build.

\section{What was Mathlib, what was new}

Mathlib: uniqueness of limits and filter congruence. Seabed: the
Stage-2 anharmonic asymptotics, consumed for the first time by the
recovery thread; the discriminant hypothesis and the
Threepoint-to-Laplace bridges arrive with them. New: only the reading
direction.

\section{Lessons}

When a thread's pattern stabilizes (this is the third
uniqueness-of-limits recovery), the marginal tide is nearly free: the
whole cost was locating the right seabed limits. The productive act was
the survey, not the proof.

\section{Follow-ups}

\begin{itemize}
\item The $\gamma$ rung: a fourth-cumulant asymptotic
($t^3 \kappa_4 \to$ explicit value) in the Stage-2 style, which would
extend both the primer programme and this recovery.
\item Tag the note's Theorem 3.1 with the recovery theorem once merged
(the note's regular case now has its first certified instance).
\end{itemize}

\end{document}
